\begin{center}
\fbox{\begin{minipage}{0.91\linewidth}
\textbf{Fresh formal verification: passed.}
The project contains 26 checked theorem declarations and 8 definitions.
The pinned Lean 4.34.0-rc2 compiler accepted the project build, both
mathematical source modules, the root import, and the 34-declaration
axiom audit. Every audited project declaration uses only
\texttt{propext}, \texttt{Classical.choice}, and \texttt{Quot.sound};
no \texttt{sorryAx} or custom axiom occurs.
A valid arithmetic control was accepted without axioms. A false arithmetic
control and an acceptance theorem missing its error-budget premise were
both rejected with exit code 1.
This certifies the finite algebra and arithmetic claims listed in this
companion, under their stated hypotheses.
\end{minipage}}
\end{center}


The focused namespace is \texttt{KomorebiQCOS}.  Its 26 theorem declarations
cover exact finite algebra and arithmetic consequences of supplied bounds.
The model supplies the real factor \(A\), uses finite column sets and finite
sample indices, and does not formalize the preprocessing that produced that
factor.  Its sample Gram convention is \(AA^\top\).  Applying a statistic
with normalization \(Q=\tfrac12v^\top Kv\) requires carrying that scalar
factor explicitly.

\begin{center}
\small
\begin{tabular}{@{}p{.25\linewidth}p{.33\linewidth}p{.33\linewidth}@{}}
\toprule
Traditional claim & Named Lean declarations & Exact formal coverage\\
\midrule
Column blocking and partition refinement &
\path{gramEntry_union}, \path{gramEntry_reindex} &
Entrywise Gram additivity for disjoint finite blocks and invariance under a
matching bijective reindexing\\[2mm]
Streamed sample-kernel action &
\path{matrixAction_gram_eq_streamed}, \path{streamedAction_union} &
Exact finite interchange of sums and block accumulation for \(Kv=A(A^\top v)\)\\[2mm]
Additivity of operator/statistic &
\path{matrixAction_gram_union}, \path{quadraticForm_gram_union} &
Disjoint block sums preserve the action and quadratic form\\[2mm]
Gram positivity &
\path{quadraticForm_gram_eq_energy}, \path{quadraticForm_gram_nonneg} &
Exact sum-of-squares identity and nonnegativity\\[2mm]
Supplied coefficient order &
\path{weightedSquares_mono}, \path{weightedSquares_tail_imp} &
Pointwise order and event implication for supplied weighted sums of squares;
no eigenvalue construction or probability measure\\[2mm]
Combining layer errors &
\path{errorBound_compose} &
Triangle-inequality addition of two assumed absolute error bounds\\[2mm]
Requested-budget acceptance &
\path{errorBound_accept} &
An assumed error bound meets an assumed larger tolerance; no machine ratio
or numerical-rounding implementation\\[2mm]
Enclosure and positivity &
\path{enclosed_positive} &
A supplied interval with a positive lower endpoint gives positivity of its
enclosed real value\\
\bottomrule
\end{tabular}
\end{center}

No row substitutes a short algebraic lemma for the full corresponding
traditional theorem.  In particular, the pointwise coefficient theorem is
one leaf of probabilistic tail sandwiching.  The enclosure theorem composes
supplied bounds; it does not prove their numerical construction.

The following remain outside this focused formal module: Gaussian rotation
and distributional reduction, the finite spectral theorem, Weyl inclusion,
phase-graph support and pruning estimates, QR spectral equivalence, residual
completion and moment inequalities, complex logarithmic packet expansion,
folding/COS inversion and analytic remainder bounds, interval-library
correctness, directed serialization, controller software, and hardware
execution.  Their evidence is given by the traditional proofs, source audits,
and numerical checks stated elsewhere.

Compiler verification and standard axiom dependencies are recorded in the
Lean companion and its declaration inventory.  The reproducible companion
includes \texttt{FiniteGram.lean}, \texttt{Certificates.lean}, and the explicit
axiom-audit entry point.  Compiler acceptance establishes the listed Lean
statements for the exact-real model; it does not establish the omitted
numerical or probabilistic claims.
